\chapter{Open obstruction classes}
\label{ch:open-obstructions}

\section{Open obstruction classes}
\label{sec:open-obstructions-and-subsequent-constructions}

The cross-level identity \(\textcircled{2}\!\leftrightarrow\!\textcircled{3}\)
has two logically distinct parts.  The Pfaffian and
orientation part,
\(\mathrm{Pf}_{\mathrm{prot}}(\mathfrak D_X^{\mathrm{DI}})=\mathcal D_X=\Delta_5\)
and \(\epsilon_o=\nu_{\Delta_5}\) on \(W^{(2)}(\La^{2,1}_{II})\),
holds on \(K3\times E\) relative to the retained D0-HN datum, retained
quotient-orientation datum, chosen Weyl lifts, and retained
\((O2_{\mathrm{atlas}})\) wall datum, together with the retained finite
geometric Pfaffian datum \((P_{\mathrm{fin}})\), of
Appendix~\ref{app:obstruction-discharges}.  The primitive recognition
\(P_X\cong\mathfrak g_{\Delta_5}\) is conditional on the finite
compact-source recognition datum \((S1)\)--\((S10)\).
The exact theorem that remains after deleting any one of these data is
Proposition~\ref{prop:ch6-hypothesis-deletion-residues}; no deletion is
filled by a scalar trace, automorphic section, or denominator product.
The implication and non-redundancy statement is
Proposition~\ref{prop:ch6-obstruction-independence}.
The finite Pfaffian packet
\(\texttt{certificates/pfaffian/k3e\_finite\_pfaffian}\) currently
records \((P_{\mathrm{fin}})\) only by its obstruction ledger
\(\texttt{blocked\_obligations.csv}\): skew perfect complexes,
Pfaffian lines, finite sections, wall charts, automorphic line
comparisons, transitions, and scalar firewalls remain missing, and the
verifier records \(\texttt{pfaffian\_certification=false}\).
The finite O2 wall-atlas packet
\(\texttt{certificates/wall\_atlas/k3e\_o2\_atlas}\) likewise records
\((O2_{\mathrm{atlas}})\) only by its obstruction ledger: simple-wall
coverage, wall objects, rank-one charts, overlaps, Weyl transports,
half-Hilbert-orbit rows, compatibility rows, type-I exclusion,
D0-limit extension, and scalar firewalls remain missing, and the
verifier records \(\texttt{o2\_certification=false}\).
The automorphic identity
\(\textcircled{1}\!\leftrightarrow\!\textcircled{2}\) is theorem
\cite{Borcherds95,GN}; the singular-theta-lift identity
\(\textcircled{2}\!\leftrightarrow\!\textcircled{5}\) is theorem
\cite{Borcherds95}; the BKM denominator computation
\(\operatorname{den}(\mathfrak g_{\Delta_5})=64^{-1}\Delta_5(2Z)\) is
theorem; the Lie-superalgebra recognition remains conditional
from the finite-stage compact \(K3\times E\) source to the closed
denominator algebra.

Each of the four obstructions sits at a definite categorical-dimensional
level on the universal chain
\(\mathsf{P}\to\mathsf{C}\to\mathsf{S}\to\mathsf{Z}\to\mathsf{A}\).
The remaining structures beyond the theorem are the Hall--Borcherds
residual comparing the level-\(n\) protected scalar shadow with a
chain-level trace datum, the chain-fusion conjecture for compact non-toric \(K3\times E\),
and the gravity-line operator algebra required to have Pentagon-face
determinant section \(\Delta_5^2\).  The Mathieu / umbral moonshine adjacency
stands between the compact source and these structures:
its connection to \(\mathfrak g_{\Delta_5}\) is conjectural.  The
conjecture asks for the twined denominator algebra and its reciprocal
partition character. The
finite protected-trace packet
\(\texttt{certificates/trace/k3e\_protected\_trace}\) presently records
the level-\(\mathsf Z\) trace datum only by its obstruction ledger:
trace categories, protected trace functors, operators, scalar
normalizations, identities, forgetful maps, gravity-residual separation,
transitions, and scalar firewalls remain missing, and the verifier
records \(\texttt{protected\_trace\_certification=false}\).
Beilinson cut separates the retained-data Pfaffian theorem from the
primitive source recognition and the gravity-line operator algebra.

The three names of \(\Delta_5\) are kept distinct on first occurrence:
the automorphic section
\(\mathcal D_X=\Delta_5\in H^0(\mathcal H_2,\mathcal L^5\otimes\nu_{\Delta_5})\)
of view \textcircled{1}; the BKM denominator
\(\operatorname{den}(\mathfrak g_{\Delta_5})=64^{-1}\Delta_5(2Z)\) of
view \textcircled{2}; and the compact protected Pfaffian
\(\mathrm{Pf}_{\mathrm{prot}}(\mathfrak D_X^{\mathrm{DI}})\) of view
\textcircled{3}, relative to the retained data of
Appendix~\ref{app:obstruction-discharges}.

\subsection{The four named obstructions}
\label{subsec:four-named-obstructions}

The four obstructions act on distinct mathematical objects before any
scalar trace is taken.  The compact target is the Dirac--Igusa pro-object
\(\mathfrak D_X^{\mathrm{DI}}=\varprojlim_R(\mathcal A_{X,R}^{E_3},
F_{X,R}^{\mathrm{hyb}},\Gamma_{X,R},\Pi_X,\Phi_R,o_R,H_R,P_R^\Pi,
C_{X,R},\Theta_{\mathrm{Kos},R},\mathcal L_{\mathrm{Pf},R},
\mathrm{pf}_{X,R},\epsilon_{o,R},\mathrm{Rec}_R)\). The level
discipline is \emph{primitive object first, cross-level identity second,
scalar shadow last}. The four obstructions and \((P_{\mathrm{fin}})\)
are pre-trace Pfaffian-orientation data of View~\textcircled{3};
the level-\(\mathsf Z\) object is the derived-centre trace pairing.
They do not identify the compact primitive Lie superalgebra; that is
the separate datum \((S1)\)--\((S10)\).

\subsubsection{(D0): D0-degeneration limit of the cosection-reduced d-critical orientation theory}
\label{subsubsec:obstruction-D0}

\paragraph{D0-degeneration.} The cosection-reduced d-critical orientation
theory on the retained Hall stack
\(\mathcal H_R^{\mathrm{or}}\) of finite HN height \(R\) extends through
the typed D0-HN datum of Definition~\ref{def:G-D0-HN-system}: a flat
D0-degeneration family, retained object/extension/flag stacks, proper
or closed HN transitions, additive K3 semiregularity cosections,
reduced vanishing-cycle specialization, Joyce--Upmeier orientation
specialization, Pfaffian-line and protected-integration compatibility,
and strict Mittag--Leffler exactness.  In particular, the orientation
line of the reduced determinant complex
\[
\mathscr D^{\mathrm{red}}_{R,\mathcal C}
=\det\operatorname{RHom}_{\mathrm{red}}(\mathcal C,\mathcal C)
\]
admits a Brav--Bussi--Dupont--Joyce--Szendr\H{o}i \cite{BBDJS2015} square-root
gerbe section over the closure
\(\overline{\mathfrak M^{\mathrm{st}}_{R,\mathcal C}}\subset
\overline{\mathcal H_R^{\mathrm{or}}}\) including the D0-degeneration
stratum, compatible with the cosection contraction inherited from
\(K3\times E\). The retained Hall datum supplies an oriented critical
finite-stage Pandharipande--Thomas/Gholampour--Sheshmani
inverse system \(\widehat{\mathcal H}^{\mathrm{or}}_{\sigma,S}
=\varprojlim_R\mathsf F_R\mathcal H_{\sigma,S}^{\mathrm{or}}\) on the
finite HN quotients of
Proposition~\ref{prop:sectorial-hall-truncation};
\textup{(D0)} asks that this system carry a coherent orientation through
the D0-degeneration limit.
The finite packet
\(\texttt{certificates/d0/k3e\_d0\_hn}\) records the missing rows in
\(\texttt{blocked\_obligations.csv}\): flat degeneration families,
retained substacks, HN transitions, derived and d-critical charts,
semiregularity cosections, vanishing-cycle and orientation
specializations, Pfaffian/protected-integration specializations,
Mittag--Leffler rows, and scalar firewalls.  Its verifier returns
\(\texttt{D0\_OBSTRUCTION\_LEDGER\_VERIFIED}\), with
\(\texttt{d0\_certification=false}\).  This is the finite D0-HN
obstruction list, not a D0-HN construction.

\paragraph{Criterion.} The retained D0-HN datum of
Definition~\ref{def:G-D0-HN-system} supplies the D0-degeneration
orientation through the retained inverse system via
Theorem~\ref{thm:G-D0}, and
Theorems~\ref{thm:G-O1}, \ref{thm:G-O1-plus}, and
\ref{thm:G-O2-orbit-transport} transport it to the type-II Weyl
character.  The Maass character \(\nu_{\Delta_5}\) and the geometric
orientation \(\epsilon_o\) then agree on
\(W^{(2)}(\La^{2,1}_{II})\).  BBDJS \cite{BBDJS2015}
vanishing-cycle data require an oriented
d-critical locus, not boundedness alone, and Joyce--Upmeier
\cite{JoyceUpmeier} square-root data on CY\(_3\) moduli do not
automatically extend to the reduced \(K3\times E\) quotient: the
cosection-induced shift in the determinant complex must be shown
compatible with descent.  The connected \(BE\) Borel/edge class
\(\alpha^{E,\mathrm{free}}=a_1u_1+a_2u_2\in H^2(BE;\mathbb F_2)\)
is separated from finite \(BE[N]\) Borel terms in the Borel
spectral sequence; the quotient-orientation descent across
object, extension, mixed, wrapped, two-step flag, overlap,
Thom--Sebastiani, and \(R'\to R\) transition strata is conditional,
with overlap/correspondence/flag/transition compatibilities named
separately.

\paragraph{Argument.} The retained Hall boundedness datum
\textup{[K3\(\times\)E-fin]}, motivated by
\cite{LiuProductStability,PiyaratneTodaModuli3folds}, supplies a
finite-type retained slice; the cited boundedness results do not by
themselves prove this slice for the compact \(K3\times E\) Hall sector.
The datum also does not give quasi-smoothness, a d-critical chart,
BBDJS coefficients, an orientation line, or compact-support
six-functor admissibility through the D0 limit.
The finite rows required for this assertion are
\(\mathrm{hn\_type\_bounds}\), \(\mathrm{semistable\_substacks}\),
\(\mathrm{derived\_enhancements}\), \(\mathrm{universal\_complexes}\),
\(\mathrm{scalar\_rigidifications}\), \(\mathrm{stratifications}\),
\(\mathrm{extension\_flag\_stacks}\), \(\mathrm{cosection\_atlas}\),
\(\mathrm{transitions}\), and \(\mathrm{scalar\_firewall}\).
The finite-moduli packet
\(\texttt{certificates/moduli/k3e\_finite\_moduli}\) presently records
the HN-bound, semistable-substack, derived-enhancement, and
rigidification-inertia rows, together with the finite class-bound
and universal-complex rows and the \(E\)-translation-rigidification
rows, retained closed-substack rows, and retained extension-closure
rows, retained HN-factor-closure rows, and retained dual-closure rows,
while its obstruction ledger keeps the flag stacks, cosection charts,
transitions, and scalar firewalls open; the verifier records
\(\texttt{moduli\_certification=false}\).
Pantev--To{\"e}n--Vaqui{\'e}--Vezzosi \cite{PTVV2013} shifted symplectic
structures cover ambient derived
moduli under Calabi--Yau hypotheses; the retained atlas is not produced
by PTVV alone. Joyce--Upmeier \cite{JoyceUpmeier} square-root orientation
on Calabi--Yau\(_3\) moduli requires exact-sequence/direct-sum
compatibility; the reduced \(K3\times E\)-quotient inherits this only
after the cosection-induced shift is shown compatible with descent.

\paragraph{Quotient-orientation classes.} The retained
quotient-orientation datum consists of
\(\alpha^{\mathrm{red}}_{R,\mathcal C}=w_2(\mathscr D^{\mathrm{red}}_{R,\mathcal C})
\in H^2(S;\mathbb F_2)\), \(\alpha^{E,\mathrm{free}}_{R,\mathcal C}\),
\(\widetilde\beta^{H}_{R,\mathcal C}\in H^2_H(S;\mathbb F_2)\), and
\(\lambda^{H}_{R,\mathcal C}\in H^1(BH;\mathbb F_2)\) on every retained
object, extension, mixed, wrapped, and flag stratum \(\mathcal C\), with
killing of the Borel mixed terms and edge-spectral-sequence
differentials at every node of the descent diagram.  On the rank-one
D6--D2--D0 stratum \(\mathcal C=(1,\beta,n)\), with \(\beta\) the
K3-fibre primitive class, the BBDJS line carries the cosection
contraction explicitly, and the four obstruction classes
\((\alpha^{\mathrm{red}},\alpha^{E,\mathrm{free}},\widetilde\beta^{H},
\lambda^{H})\) are the retained quotient-orientation datum of
Appendix~\ref{app:G-O1-discharge}.  The remaining source-side problems
are this finite-stabilizer quotient orientation and the finite
compact-source recognition datum for the Hall primitive algebra.
The finite orientation packet
\(\texttt{certificates/orientation/k3e\_reduced\_orientation}\) records
this as an obstruction ledger, not as an orientation proof.  Its
\(\texttt{blocked\_obligations.csv}\) table lists the missing
square-root, Thom--Sebastiani, quotient Borel, finite-stabilizer,
Weyl-lift, Coxeter, transition, and scalar-firewall rows, and the
verifier reports \(\texttt{ORIENTATION\_OBSTRUCTION\_LEDGER\_VERIFIED}\)
with \(\texttt{orientation\_certification=false}\).

\subsubsection{(O1): strong reduced orientation on the K3{\(\times\)}E-quotient self-Ext frame bundle}
\label{subsubsec:obstruction-O1}

\paragraph{Reduced quotient orientation.} The BBDJS \cite{BBDJS2015} orientation gerbe on the moduli
of stable sheaves on \(K3\times E\) descends to the
\(K3\times E\)-quotient self-Ext frame bundle, with reduced orientation
gauge: on the stable stratum
\(\mathfrak M^{\mathrm{st}}_{R,\mathcal C}\subset\mathcal H_R^{\mathrm{or}}\)
the obstruction tuple
\[
\mathfrak o^{\mathrm{or}}_{R,\mathcal C}
=\bigl(\alpha^{\mathrm{red}}_{R,\mathcal C},\,
\alpha^{E,\mathrm{free}}_{R,\mathcal C},\,
\{(\widetilde\beta^{H}_{R,\mathcal C},\lambda^{H}_{R,\mathcal C})\}_{H},\,
\mu^{\mathrm{or}}_{R,e},\,\dots\bigr)
\]
vanishes simultaneously, with all overlap, correspondence, flag, and
\(R'\to R\) transition data trivialized coherently. \textup{(O1)}
demands the strong reduced descent: the connected-\(BE\) class
\(\alpha^{E,\mathrm{free}}\), the finite-\(BE[N]\) gerbe classes
\(\widetilde\beta^H\), and the residual finite linearization characters
\(\lambda^H\) all vanish, separately, on every retained stratum, and the
spectral-sequence differentials and Borel mixed terms are killed.

\paragraph{Criterion.} The Hall well-formedness criterion is prior to
recognition: by
Proposition~\ref{prop:hall-product-only-no-primitives}, an associative
Hall product without a counital coproduct, bialgebra compatibility, Hopf
pairing, radical, quotient, and strict transition data has no defined
primitive source object.  The primitive recognition criterion then reads:
target arithmetic on \(W_{\le3}\) is target
Borcherds--Gritsenko--Nikulin--Kac data only, and recognition
requires an actual compact source datum supplying
\(M,D,B,G,K,Q,A_\beta\), relation rows, Hopf radical/coideal kernel
equality, PBW associated-graded comparison, and strict
transition/Mittag--Leffler rows.  The Koszul comparison adds the finite
comparison residual
\(\mathfrak O_{\mathrm{Kcmp},R}\): by
Proposition~\ref{prop:primitive-restriction-not-koszul}, agreement of
the source-to-target map on constant-current primitives does not imply
the quasi-isomorphism
\(\Omega_E^{\mathrm{ch}}C_X\simeq\mathsf{Den}_{\Delta_5,E}\), and by
Definition~\ref{def:finite-koszul-comparison-datum} even that
quasi-isomorphism must still respect Weyl transport, Pfaffian
orientation, Hall pairing, radical quotient, PBW filtration, and HN
transitions.  The pro-object comparison also requires the finite
morphism rows
\(\mathrm{cofinal\_subsystems}\), \(\mathrm{stage\_objects}\),
\(\mathrm{stage\_morphisms}\), \(\mathrm{component\_maps}\),
\(\mathrm{composition\_laws}\), \(\mathrm{ml\_exactness}\),
\(\mathrm{pro\_isomorphisms}\), and \(\mathrm{scalar\_firewall}\);
without them a scalar Pfaffian product or denominator identity does not
define a Dirac--Igusa pro-object.
Theorem~\ref{thm:G-O1} of
Appendix~\ref{app:obstruction-discharges} states what the retained
\textup{(O1)} datum carries
to the recognition theorem. The connected \(E\) and finite \(E[N]\)
cohomologies are \emph{distinct} obstruction spaces:
\(H^*(BE;\mathbb F_2)=\mathbb F_2[u_1,u_2]\) with
\(|u_i|=2\), so \(H^1(BE;\mathbb F_2)=0\); for
\(E[2]\cong(\mathbb Z/2)^2\),
\(H^*(BE[2];\mathbb F_2)=\mathbb F_2[x_1,x_2]\) with
\(|x_i|=1\), and the residual linearization
\(\lambda=\lambda_1x_1+\lambda_2x_2\) is independent data after the
degree-two gerbe class is killed.

\paragraph{Argument.} Translation invariance under \(E\)-action does
not produce orientation descent: a fixed underlying line may carry a
nontrivial finite character, and ordinary translation invariance is
\emph{not} a substitute for vanishing of
\(\lambda^H\) on full two-primary generators of
\(H^*(B(\mathbb Z/2^a)^2;\mathbb F_2)
=\Lambda(x_1,x_2)\otimes\mathbb F_2[y_1,y_2]\) at \(N=2^a\). The
quotient-orientation criterion is finite \v{C}ech--Borel descent on the
retained groupoid, not a vanishing theorem.

\paragraph{Quotient-orientation theorem.} Theorem~\ref{thm:G-O1} is the
quotient-orientation criterion relative to the retained
\((O1_{\mathrm{quot}})\) datum.  That datum contains the BBDJS gerbe
sections on the self-Ext frame bundle, their null classes on the
rank-one D6--D2--D0 stratum and its extension, mixed, wrapped, and
two-step flag completions, and the overlap and correspondence
transitions on the cofinal HN system.  The data
\(M,D,B,G,K,Q,A_\beta\), the relation rows, the Hopf radical equality,
and the PBW comparison required by primitive recognition are still
separate compact-source data.

\subsubsection{(O1)\(^+\): Weyl-equivariant transport of the orientation along W\((^2)\)(\(\La_{II}^{2,1}\)) reflections}
\label{subsubsec:obstruction-O1plus}

\paragraph{Weyl transport.} The orientation character on the type-II Weyl
group \(W^{(2)}(\La^{2,1}_{II})\) is Weyl-equivariant: the orientation
line of \textup{(O1)} is equipped with coherent lifts
\(\tau_i:o_{\mathcal C}\to s_{\delta_i}^{*}o_{\mathcal C}\) of the
three type-II simple reflections, each acting with sign \(-1\), and
the projective determinant-line cocycle
\(c_o\in Z^2(W^{(2)}(\La^{2,1}_{II});\mathbb F_2)\) of the lift is a
coboundary, with quotient torsor defects
\(\omega_{i,\mathcal C}=0\) on every retained groupoid component. At
finite height \(R\), the partial action groupoid \(\mathcal G_R\)
carries
\(c_{o,R}\in Z^2(\mathcal G_R;\underline{\mathbb F}_2)\), and
\textup{(O1)\(^+\)} asks that
\([c_{o,R}]=0\), \(\tau_i^2=1\), and \(\omega_{i,\mathcal C}=0\) hold
coherently on every stratum and at every transition height.

\paragraph{Criterion.} The local sign-extraction criterion reads:
only after the unit character is killed and
\(N_\delta^{\mathrm{Pf}}=1\) is computed from the reduced normal
self-Ext chart does the local sign equal \(-1\), and Maass values
\(\nu_{\Delta_5}(s_{\delta_i})=-1\) plus divisor orders
\(d^{\mathrm{aut}}_{\delta_i}=f(0,\pm1)=1\) supply target data
only --- they do not compute \(N_\delta^{\mathrm{Pf}}\),
\(\chi_\upsilon\), or \(\epsilon_o\). The Coxeter presentation
\[
W^{(2)}(\La^{2,1}_{II})
=\langle s_{\delta_1},s_{\delta_2},s_{\delta_3}\mid s_{\delta_i}^2=1\rangle
\]
implies that any orientation character on \(W^{(2)}(\La^{2,1}_{II})\)
is determined by three generator signs once \textup{(O1)\(^+\)} gives
the honest Weyl lift; the three rank-one signs established by
\((O2_{\mathrm{atlas}})\) are \(-1\), and then propagate to the full
character.  The
Weyl-equivariant transport is a separate identity, distinct from the
no-extra-relations and PBW identities, established relative to the Weyl lifts in
Theorem~\ref{thm:G-O1-plus}.

\paragraph{Argument.} Three local signs on the simple reflections
imply a character on the full type-II Weyl group only after
\([c_{o,R}]=0\) and \(\omega_{i,\mathcal C}=0\): the Weyl
determinant-line cocycle measures whether the lift is honest, and the
quotient torsor defects measure whether the descent commutes with the
reflection action on the retained groupoid. The cocycle reduces to the
classical group-cohomology class only on a complete
\(W^{(2)}(\La^{2,1}_{II})\)-orbit with constant coefficients; finite
height \(R\) is a partial action groupoid whose component data must be
identified on intersections.

\paragraph{Conditional theorem.} Theorem~\ref{thm:G-O1-plus} is the
criterion for finite-height lifts \(\tau_{i,R}\) whose squares are
one, whose projective cocycle class \([c_{o,R}]\) is zero, and whose
torsor defects \(\omega_{i,\mathcal C}\) vanish on the retained
groupoid components.  These data give Weyl-equivariant
orientation transport once supplied.  In the current finite
orientation packet they are exactly the missing Weyl-lift and
Coxeter rows in the obstruction ledger.  They do not supply the
height-seven source matrices required for the relation-closed
primitive recognition window.

\subsubsection{(O2): local Pfaffian wall normal form on type-II walls}
\label{subsubsec:obstruction-O2}

\paragraph{Type-II wall normal form.} On the three type-II walls
\(\Hpl_{\delta_i}\) for
\(\delta_1=2f_2-f_3,\,\delta_2=2f_{-2}-f_3,\,\delta_3=f_3\), the
retained wall-atlas datum has half-Hilbert sign support of size
\(4096=2^{12}\).  The local computation in Appendix~E is only the
rank-one crossing. Concretely, in an
\(s_\delta\)-equivariant real-analytic d-critical/Kuranishi chart
\((U_\delta,x_\delta,u_\delta,K^{\mathrm{red}}_\delta)\), with
\(K^{\mathrm{red}}_\delta\) the real form of the cosection-reduced
self-Ext complex
\(\operatorname{RHom}_{\mathrm{red}}(\mathcal C_\delta,\mathcal C_\delta)\),
the splitting
\(K^{\mathrm{red}}_\delta\simeq K^{\mathrm{tan}}_\delta\oplus
K^{\mathrm{nor}}_\delta\) is equipped with the extension class
\(\eta_\delta\in
\operatorname{Ext}^1_{D^{s_\delta}(U_\delta)}
(K^{\mathrm{nor}}_\delta,K^{\mathrm{tan}}_\delta)\) trivial after
the reduced \(E\)-quotient and quotient-orientation trivializations,
and the normal Pfaffian admits the rank-one form
\(\operatorname{Pf}(K^{\mathrm{nor}}_\delta)
=\upsilon_\delta\,u_\delta^{N_\delta^{\mathrm{Pf}}}\) with
\(\upsilon_\delta(x_\delta)\ne0\),
\(s_\delta(\upsilon_\delta)=\upsilon_\delta\), and
\(N_\delta^{\mathrm{Pf}}=1\).

\paragraph{The 4096 sign sum.} The integer \(4096=2^{12}\) is the
half-Hilbert sign count retained in \((O2_{\mathrm{atlas}})\): twelve
sign degrees of freedom, each \(\mathbb Z/2\), give \(2^{12}=4096\)
signed states.  The type-II wall locus realizes this count only after
the component, boundary, overlap, quotient, and protected-integration
compatibilities of the retained atlas hold. The arithmetic
match \(-4096=-(64)^2\) with the OP/DT scalar
\(Z^X_{\mathrm{OP}}=-4096\,\Delta_5^{-2}\) is the squared-theta
normalization, fixed by Borcherds normalization and the OP scalar
sign convention; it is \emph{not} a derivation of \(\epsilon_o\).

\paragraph{Criterion.} The local model and the exact separation
from OP and Maass scalars are the following: scalar
constants \(D_5=64^{-1}\Delta_5\),
\(\chi_{10}^{\mathrm{OP}}=4096^{-1}\Delta_5^2\),
\(Z^X_{\mathrm{OP}}=-4096\,\Delta_5^{-2}\) are scalar normalizations
and do not define \(\epsilon_o\); the Pfaffian normal form is the
content.  Proposition~\ref{prop:scalar-pfaffian-data-not-section}
strengthens this separation: the scalar product, its square, and the OP
trace also do not determine \(\operatorname{pf}_{X,R}\) as a section of
the Pfaffian line, because the order-two monodromy is invisible after
squaring.  Theorem~\ref{thm:G-O2-orbit-transport} is relative to
\((O2_{\mathrm{atlas}})\), whose matrices over \(\mathbb Z\) or
\(\mathbb Q\) record the Pfaffian-rank identity
\(N_\delta^{\mathrm{Pf}}=1\), the splitting \(\eta_\delta=0\), the
unit \(\upsilon_\delta\) with \(s_\delta\)-character trivial, and
overlap/transition compatibility across the cofinal HN system.
The packet
\(\texttt{certificates/wall\_atlas/k3e\_o2\_atlas}\) presently supplies
only the obstruction ledger
\(\texttt{O2\_OBSTRUCTION\_LEDGER\_VERIFIED}\), with
\(\texttt{o2\_certification=false}\).  It records the missing atlas
data; it does not supply the wall atlas.

\paragraph{Argument.} The condition \((\delta,\delta)=2\), the
automorphic divisor order
\(d^{\mathrm{aut}}_{\delta_i}=1\), and the Maass character
\(\nu_{\Delta_5}(s_{\delta_i})=-1\) are target data that do not imply
the Pfaffian rank assertion
\(N_\delta^{\mathrm{Pf}}=1\). Graph-isogeny atom candidates
\(B_2=i_{\Gamma_\varphi,*}L\) are Koszul-Ext-local; they are
\emph{candidates}, not constructions, until semistability, full
\(\widehat\Gamma_X\) charge match, reduced normal Ext splitting,
invariant unit, and atlas overlap compatibility are established. The
``higher terms do not alter'' phrase in the rank-one Ext criterion
requires the explicit equivariant real/parametric Morse lemma and
unit check before it operates as a theorem.

\paragraph{Wall-atlas transport.}
Definition~\ref{def:G-O2-atlas} gives the retained type-II wall-atlas
datum.  Theorem~\ref{thm:G-O2-orbit-transport} transports that datum
over the half-Hilbert orbit; it does not construct the graph-isogeny
or reducible-curve atoms.  Theorem~\ref{thm:G-R1-R2} identifies the
\(2^{12}\) sign count with the squared theta-leading normalization
only after scalarization.
The remaining relation-closed source comparison is the recognition
datum \((S1)\)--\((S10)\), not a type-II wall normal-form problem.

\subsection{Mathieu and umbral moonshine adjacency}
\label{subsec:moonshine-adjacency}

The K3 elliptic genus has a precise representation-theoretic
decomposition. Eguchi, Ooguri, and Tachikawa \cite{EOT} observed in 2010 that
the elliptic genus
\(Z_{\mathrm{K3}}(\tau,z)=2\phi_{0,1}(\tau,z)\) admits an
\(M_{24}\)-graded character expansion: the mock-modular completion
\(\Sigma(\tau)\) of the elliptic genus is a graded \(M_{24}\)-module,
and the multiplicities of irreducible \(M_{24}\)-representations in
the first several Fourier coefficients of the holomorphic part match
the dimensions of irreducible \(M_{24}\)-representations on the nose.
Gannon \cite{GannonUmbral} proves the arithmetic identity; the
geometric origin via a sigma-model action of \(M_{24}\) on the K3
elliptic-genus state space remains open outside the classes accounted
for by known K3 sigma-model symmetries.  The \(M_{24}\) action
visible at the level of the elliptic genus does not lift to a global
geometric symmetry of \(D^b(\mathrm{Coh}(K3))\); the work of
Gaberdiel--Hohenegger--Volpato \cite{GHV} on Mathieu twining characters
identifies the precise refinement at the level of supersymmetric
sigma-model conformal field theory data.

The umbral moonshine of Cheng--Duncan--Harvey \cite{CDH} extends Mathieu
moonshine to a family of \(23\) cases indexed by Niemeier root
systems, one case per Niemeier lattice having a non-empty root system.
For a Niemeier lattice \(N^X\) with root system \(X\), the finite
group is the umbral group
\[
G^X=\operatorname{Aut}(N^X)/W^X,
\]
not the discriminant group.  Mathieu moonshine is the \(X=A_1^{24}\)
case.  The Leech lattice is the \(24\)th Niemeier lattice and has empty
root system; it belongs to Conway moonshine, not to the \(23\)-case
umbral list.  The umbral moonshine conjecture of
Cheng--Duncan--Harvey \cite{CDH}, in the module-existence form proved by
Gannon \cite{GannonUmbral}, asserts that each pair \((X,G^X)\) carries a
graded module whose graded character is the prescribed vector-valued
mock modular form.

\paragraph{The conjectural connection to \(\mathfrak g_{\Delta_5}\).}
The Mathieu moonshine module lies on the \(N=4\) superconformal
character side of the K3 elliptic genus.  The Eguchi--Ooguri--Tachikawa
module is graded by the massive level \(n\); the \(z\)-dependence and
the \(l\)-index are carried by the \(N=4\) character, not by a primitive
\((n,l)\)-grading on the \(M_{24}\)-module.  The Beilinson-cut
statement is therefore a conjectural extension of the twined K3
elliptic genus to the Borcherds--Kac--Moody datum on
\(\La^{2,1}_{II}\), producing a \(g\)-twined refinement of
\(\mathfrak g_{\Delta_5}\) with denominator
\(\operatorname{den}(\mathfrak g_{\Delta_5}^{(g)})\) for each
\(M_{24}\)-conjugacy class \(g\). The umbral cases of
\cite{CDH,ChengHarrison} extend the conjecture to a family of \(23\)
Borcherds--Kac--Moody-style algebras associated to the \(23\) Niemeier
root systems with non-empty root sets.

\paragraph{Precise conjecture.}
For each conjugacy class \(g\in M_{24}\), let
\(\phi^{(g)}_{0,1}(\tau,z)\) denote the \(g\)-twined K3 elliptic
genus of \cite{GHV}.  The Cheng--Duncan--Harvey--Gannon umbral
moonshine conjecture, refined to the BKM stratum of
\(\La^{2,1}_{II}\), asserts the existence of a \(g\)-twined
Borcherds--Kac--Moody algebra
\(\mathfrak g_{\Delta_5}^{(g)}\) with denominator
\[
 \operatorname{den}\bigl(\mathfrak g_{\Delta_5}^{(g)}\bigr)
 \;=\;\Theta\bigl(\phi^{(g),\mathrm{Weil}}_{0,1}\bigr)\!\bigm|_{\Hpl_2},
\]
the singular theta lift of the \(g\)-twined Weil-representation
extension of \(\phi^{(g)}_{0,1}\), with the chamber and Weyl-vector
data determined by the relevant umbral pair \((X,G^X)\), when such a
pair is part of the twining datum.  No map from all \(M_{24}\)-classes
to Niemeier root systems is asserted here.

\paragraph{The order-seven and order-fifteen classes.}
The four classes \(\{7A,7B,15A,15B\}\) lie outside the diagonal
Gritsenko--Cl\'ery/Borcherds frame used here, namely the rows selected
by \(\varphi(N)\mid2\); this does not exclude their occurrence in
other CHL dyon towers.  Their cyclotomic fields have degrees \(6\) and
\(8\).  Their cycle shapes on the standard \(24\)-dimensional permutation
representation are
\[
 7A:\ 1^3\!\cdot 7^3,\qquad
 7B:\ 1^3\!\cdot 7^3,\qquad
 15A:\ 1\!\cdot 3\!\cdot 5\!\cdot 15,\qquad
 15B:\ 1\!\cdot 3\!\cdot 5\!\cdot 15,
\]
the \(7A\)/\(7B\) pair distinguished by Galois conjugation under
\(\mathrm{Gal}(\Q(\zeta_7)/\Q)\) and the \(15A\)/\(15B\) pair by
\(\mathrm{Gal}(\Q(\zeta_{15})/\Q)\); the rational character on each
pair is integer-valued, but individual class characters take values
in the cyclotomic field of order equal to the class order.  These
cycle shapes have different geometric status: \(7A\) and \(7B\) have
six orbits and pass the Mukai--Kondo orbit test, while \(15A\) and
\(15B\) have four orbits and are non-geometric in that test
\cite[\S2]{GaberdielVolpatoOrbifoldK3}.  The orders
\(|g|\in\{7,15\}\) satisfy \(\varphi(|g|)\nmid 2\)
(\(\varphi(7)=6\), \(\varphi(15)=8\)), placing them outside the
Gritsenko--Cl\'ery frame \(N\in\{1,2,3,4,6\}\) of \(\kappa_{\mathrm{BKM}}\)-universality
selected by \(\varphi(N)\mid 2\) of
Theorem~\ref{thm:bkm-kappa-universal}.  A Borcherds--Kac--Moody
construction for these twined genera therefore requires data beyond
that denominator frame.

\begin{conjecture}[Order-seven and order-fifteen twined denominators]
\label{conj:order-seven-fifteen-twined-denominators}
For each \(g\in\{7A,7B,15A,15B\}\):
\begin{enumerate}
\item[\textnormal{(M${}^\star$\!7A)}] The Gaberdiel--Hohenegger--Volpato
\(7A\)-twined K3 elliptic genus
\(\phi^{(7A)}_{0,1}\) is a weight-\(0\), index-\(1\) weak Jacobi form
with multiplier for the appropriate congruence subgroup attached to
the \(7A\) class.  Its singular theta lift, if the Weil extension and
Weyl chamber data are fixed, has Borcherds weight
\(c_{7A}(0)/2\) and multiplier, and equals
\(\operatorname{den}(\mathfrak g_{\Delta_5}^{(7A)})\) for a
conjectural \(7A\)-twined Borcherds--Kac--Moody algebra.
\item[\textnormal{(M${}^\star$\!7B)}] Identical statement for the
Galois conjugate class \(7B\), with
\(\phi^{(7B)}_{0,1}=\phi^{(7A)}_{0,1}\circ
\sigma_{\zeta_7\to\zeta_7^3}\) the Galois-conjugate character.
\item[\textnormal{(M${}^\star$\!15A)}] The
\(15A\)-twined character \(\phi^{(15A)}_{0,1}\) is a weight-\(0\),
index-\(1\) weak Jacobi form with multiplier for the appropriate
congruence subgroup attached to the \(15A\) class.  Its singular theta
lift, if the Weil extension and Weyl chamber data are fixed, has
Borcherds weight \(c_{15A}(0)/2\) and multiplier, and equals
\(\operatorname{den}(\mathfrak g_{\Delta_5}^{(15A)})\).
\item[\textnormal{(M${}^\star$\!15B)}] Identical statement for the
Galois conjugate class \(15B\).
\end{enumerate}
The four clauses have finite necessary conditions in their first
Fourier coefficients, but no finite truncation proves the denominator
identity.
\end{conjecture}

\begin{remark}[Fourier conditions and the parity obstruction]
\label{rem:order-seven-fifteen-cross-checks}
Three Fourier conditions constrain each clause
\((\mathrm M^\star\!g)\) at finite truncation depth.

\textup{(i)} The leading Fourier coefficient
\(c_g(0)=[q^0 r^0]\phi^{(g)}_{0,1}\) is the
Gaberdiel--Hohenegger--Volpato McKay--Thompson value of \(g\) at the
\(q^0 r^0\) position of the K3 elliptic genus, tabulated in
\cite[Table~1]{GHV} for all 26 conjugacy classes.  Outside the CHL
frame \(\{1A,2A,3A,4A,6A\}\), this leading coefficient need not be an
even integer in the Borcherds lattice normalisation; for the
classes \(\{7A,7B,15A,15B\}\) it is not one of the values
\(c_N(0)\in\{10,6,4,3,2\}\) of
Theorem~\ref{thm:bkm-kappa-universal} \textup{(within the frame the
\(N=4\) entry \(c_4(0)=3\) is already odd, with half-integral weight
\(\tfrac32\))}.

\textup{(ii)} The Borcherds weight formula gives
\(\kappa_{\mathrm{BKM}}(\Phi^{(g)})=c_g(0)/2\) only after the twined
Borcherds datum and its lattice have been fixed.  For the
order-seven and order-fifteen classes this weight need not be an integer; the lift must
therefore be treated as a multiplier-bearing automorphic product, not
as a weight-\(5\) paramodular form.  This is the arithmetic obstruction
to extending the level-\(\mathsf Z\) denominator construction beyond
the CHL frame.

\textup{(iii)} Mathieu-moonshine consistency with the EOT
\(M_{24}\)-multiplicity predictions through the first eight Fourier
coefficients of \(\Sigma(\tau)\) imposes a finite necessary condition
on each clause; proving
\((\mathrm M^\star\!g)\) requires identifying the geometric origin
of the \(M_{24}\)-action on the K3 elliptic-genus state space and its
umbral refinement in the sense of Cheng--Harrison \cite{ChengHarrison}.
\end{remark}

\paragraph{Independence from the four Pfaffian--Dirac obstructions.}
The Pfaffian theorem~\ref{thm:ch6-pfaffian-dirac} and the
orientation-character theorem~\ref{thm:ch6-orientation-character} are
retained-data theorems on \(K3\times E\), after the retained D0-HN datum,
retained quotient-orientation datum, chosen Weyl lifts,
\((O2_{\mathrm{atlas}})\), and the retained finite geometric Pfaffian
datum \((P_{\mathrm{fin}})\) of
Theorem~\ref{thm:G-four-obstruction-criterion}.  The primitive
recognition theorem remains conditional on \((S1)\)--\((S10)\), and
the level-\(\mathsf A\) gravity-line operator algebra remains conditional on
the Vol~II Hall--Borcherds residual.  The moonshine
adjacency is independent of both: it lives at level \(\mathsf{Z}\)
(centre/denominator) and at level \(\mathsf{A}\) (acting object on
the K3 elliptic-genus state space).  A proven identification of the
moonshine module as a \(g\)-twined character of
\(\mathfrak g_{\Delta_5}\) does not by itself solve the Vol~II
Hall--Borcherds residual; conversely the Hall--Borcherds residual
does not prove the moonshine adjacency.

\paragraph{Twined denominator problem.} The \(g\)-twining operation would
produce \(\operatorname{den}(\mathfrak g_{\Delta_5}^{(g)})\) from
\(\operatorname{den}(\mathfrak g_{\Delta_5})\) and an \(M_{24}\)-action
on a denominator-compatible lift of the Mathieu module.  The ordinary
massive-level grading of the \(M_{24}\)-module supplies the mock modular
coefficients; the Borcherds product requires the additional Jacobi
variable and \(\La^{2,1}_{II}\)-root grading.  That lift, and the
denominator-compatible action on it, are not part of the EOT theorem.

\paragraph{Compatibility with Vol III.} The Vol~III dimension-stratified
BKM catalogue \cite{CYQGVolIII} places \(\mathfrak g_{\Delta_5}\) as
the \(d=3\) sibling, with companions at \(d=3\) (Borcherds Monster
\(V^\natural\) on \(\mathrm{II}_{1,1}\)) and \(d=5\) (Fake Monster on
\(\mathrm{II}_{25,1}\) via \(K3\times K3\times E\)).
The Mathieu/umbral moonshine adjacency to \(\mathfrak g_{\Delta_5}\)
is a level-\(\mathsf{Z}\) cross-stratum conjecture: the denominator-compatible
\(M_{24}\) action on the K3 elliptic-genus state space for
\(\{7A,7B,15A,15B\}\) is the missing piece on the K3-fibre side; the
\(g\)-twined denominator
\(\operatorname{den}(\mathfrak g_{\Delta_5}^{(g)})\) is the missing
piece on the BKM side.

\paragraph{Twined denominator identity.} For each of the \(26\)
conjugacy classes of \(M_{24}\), the expected object is a
\(g\)-twined Borcherds product on \(\La^{2,1}_{II}\).  Its
specialization to the \(N=4\) character expansion must recover the
twined McKay--Thompson coefficients of EOT--Gannon.  A coefficientwise
refinement of the EOT multiplicities by BKM denominator coefficients is
therefore a consequence only after the denominator-compatible Jacobi
lift and the \(\La^{2,1}_{II}\)-root grading have been fixed.  The
required identity is the consistency of the \(g\)-twined
denominator with the OP scalar
\(Z^X_{\mathrm{OP}}=-4096\,\Delta_5^{-2}\) under the \(g\)-twining
construction, including the squared sign convention.

\subsection{Residual operator structures}
\label{subsec:subsequent-constructions}

The Beilinson cut separates the remaining operator data by the object
on which the missing structure must live: the Hall--Borcherds operator
lift, the compact non-toric chain-fusion equivalence, and the
gravity-line algebra.

\subsubsection{The Hall--Borcherds residual}
\label{subsubsec:open-hall-borcherds}

\paragraph{Hall--Borcherds residual.} The level-\(n\) protected scalar shadow
\(Z^{K3\times E}_{\mathrm{BPS}}=(\Phi_{10}^{\mathrm{un}})^{-1}=
\Delta_5^{-2}\) on the closed
\(K3\times E\mapsto\,\mathrm{point}\) cobordism is the BPS partition
function.  A 3d gravitational path-integral interpretation with a
genuine operator-level lift requires the Hall--Borcherds residual: an
operator-valued completion of the scalar trace whose trace recovers the
level-\(n\) shadow and whose algebraic content carries the metric
degrees of freedom.  This residual remains open in Vol~II.

\paragraph{Comparison with Vol II.} Vol~II \cite{ChiralBarCobarVolII}
names the gravity-line operator algebra required to have Pentagon-face
determinant section \(\Phi_{10}^{\mathrm{un}}=\Delta_5^2\), whose
partition character is \((\Phi_{10}^{\mathrm{un}})^{-1}=\Delta_5^{-2}\), and identifies the
Hall--Borcherds residual as the comparison arrow from the bulk
\(Z^{\mathrm{der}}_{\mathrm{ch}}(A_b)\) to the line-operator algebra.
The Vol~II residual asks for the construction of the
gravity-line operator algebra with Pentagon-face determinant
\(\det_{\mathrm{Pentagon}}=\Phi_{10}^{\mathrm{un}}=\Delta_5^2\)
on the closed cobordism, whose primitive square-root Borcherds weight is
\(\omega_{\mathrm{Borcherds}}(\Delta_5)=c_1(0)/2=5\) at the
K3-Heisenberg witness.
The Igusa-side datum is the same automorphic determinant identity
\(\Delta_5^2=\Phi_{10}^{\mathrm{un}}\), with the Appendix~G Pfaffian datum
\(\mathfrak D_X^{\mathrm{DI}}\) carrying the retained orientation and
Weyl-transport data on the chiral side.

\paragraph{Argument.} The protected graded trace
\[
\mathrm{Tr}_{\mathrm{prot}}\bigl((-1)^F
q^{L_0}r^{J_0}s^{\widetilde{L}_0}\bigr)=\Delta_5^{-2}
\]
forgets orientation; the OP chamber representative is
\(-4096\) times this trace.  The
sign-character \(\epsilon_o\) is invisible in the squared form and is
part of the Pfaffian--Dirac orientation theorem, not of the
gravity-line residual.  The Hall--Borcherds residual asks instead for
an acting operator algebra whose trace recovers the scalar shadow.
The realization as a three-dimensional gravitational path integral is conditional on
Brown--Henneaux, modular invariance, vacuum dominance, and saddle
dominance \cite{HM}; no such realization is asserted.

\paragraph{Scalar non-faithfulness.} The structural reason
\(\Delta_5^{-2}\) does not by itself recover the operator-level
gravity-line algebra is the scalar non-faithfulness of the graded
Euler character on filtered \(SC^{\mathrm{ch},\mathrm{top}}\)-objects,
proved as a lemma in Vol~III~\cite{CYQGVolIII}: an equality
\(\chi(C_1) = \chi(C_2)\) between filtered
\(SC^{\mathrm{ch},\mathrm{top}}\)-objects does not determine a
quasi-isomorphism, an \(SC^{\mathrm{ch},\mathrm{top}}\)-morphism, a
Hall product, a Drinfeld pairing, a current-envelope locality
structure, or a derived-centre trace --- adding a filtered acyclic
summand \(K\) with \(\chi(K) = 0\) leaves the scalar unchanged while
admitting non-trivial extensions of the operadic action, and a
filtered differential may be perturbed within the associated graded
without altering \(\chi\). The missing chain-level content underneath
\(\Delta_5^{-2}\) sits under Vol~III's Hall--Borcherds datum: a
positive-half Hall--Borcherds bialgebra morphism on \(E\), a completed
Drinfeld-double extension, a current-envelope morphism along \(E\), and
trace compatibility with \(\Delta_5^{-2}\) \cite{CYQGVolIII}.  The
Igusa scalar is only the trace shadow of this datum.  The
Hall--Borcherds residual is the missing chain-level lift
together with its finite-window identities.

\paragraph{Pentagon-face cyclic trace.} The required cyclic trace is
\(\mathrm{Tr}_{\mathrm{Pentagon}}^{\mathrm{cl}}\) on the closed-colour
factorization algebra \(F^{\mathrm{cl}}\) at the
\(K3\times E\)-Heisenberg witness.  Its evaluation must recover
\(c_1(0)/2=5\) on the universal Borcherds-weight identity, and its
operator-level lift to the open-colour boundary algebra
\(A_b=\mathrm{End}_{\mathcal C}(b)\) must factor through the
chain-fusion datum below.  The compatibility condition is the
consistency of the Pentagon-face scalar with the unscaled trace
\(\Delta_5^{-2}\) under the squared-trace operation.  The OP chamber
representative is \(-4096\) times this trace.

\subsubsection{Chain fusion for compact non-toric K3{\(\times\)}E}
\label{subsubsec:open-chain-fusion}

\paragraph{Chain-fusion datum.} For a Calabi--Yau\(_d\) category
\(\mathcal C_X\) and a chart datum \((\Sigma_{d-1},C)\), a
chain-fusion datum consists of a boundary vacuum
\(b(X,\Sigma,C)\) in an open factorization dg-category on
\((C,D_C,\tau_C)\), a comparison morphism
\[
 \Phi^{(\Sigma_{d-1},C)}_d(\mathcal C_X)
 \longrightarrow
 \mathrm{End}_{\mathcal C}(b(X,\Sigma,C)),
\]
and the compatibility of this morphism with clutching, trace, and the
Calabi--Yau orientation.  The chain-fusion conjecture asserts that this
morphism is a quasi-isomorphism for the admissible compact chart.  The
affine toric examples \(\mathbb C^3\), local \(\mathbb P^2\), and the
conifold furnish local comparison models; they do not construct the
compact non-toric boundary vacuum for \(K3\times E\).  At \(d=3\), the
retained Dirac--Igusa data supply the protected scalar shadow on
\(K3\times E\); the boundary-vacuum construction required for chain
fusion remains open.

\paragraph{Comparison with Vol III.} Vol~III's catalogue
\cite{CYQGVolIII} names this conjecture in the two-stage
factorization frame and identifies, after chain fusion, the
level-\(\mathsf Z\) Drinfeld double obligation
\(G(X)=D(Y^+(X))\) for compact non-toric \(\mathrm{CY}_3\), together
with the conditional TCFT/framing/anomaly/completion datum needed
for compact non-formal \(\mathrm{CY}_3\) in general. Chain fusion
remains a conjecture, not a theorem; the assertion of \(\Delta_5\) as
a protected Pfaffian section on
\(K3\times E\) is the retained-data identity
\(\mathrm{Pf}_{\mathrm{prot}}(\mathfrak D_X^{\mathrm{DI}})=\mathcal D_X=\Delta_5\)
of view \textcircled{3}, not an underlying microscopic state space.

\paragraph{Argument.} The compact \(K3\times E\) case at \(d=3\) is a
candidate instance of chain fusion via the rectified finite
\(K3\)-to-\(E\) specialization
(Definition~\ref{def:O1-strong-reduced-orientation}).
The proposed chiral specialization, if constructed from the retained
finite-stage datum, has target form
\(\mathrm{Sp}^{\mathrm{ch}}_{K3,E}(\mathfrak A_{K3\times E})
=U^{\mathrm{ch}}(\mathfrak{heis}_{\mathrm{Muk}})\otimes
U^{\mathrm{ch}}(\mathfrak g^{\mathrm{BPS}}_{K3})\)
on \(\mathrm{Ran}(E)\).  The map
\(\mathrm{Sp}^{\mathrm{ch}}_{K3,E}\) is the chosen
\([K3{\to}E\text{-}\mathrm{sp}]\) datum, not a canonical functor.  The
finite proof rows for that datum are
\(\mathrm{formal\_target\_charts}\), \(\mathrm{field\_complexes}\),
\(\mathrm{e3\_operations}\), \(\mathrm{bv\_qme}\),
\(\mathrm{anomaly\_framing}\), \(\mathrm{compact\_support}\),
\(\mathrm{factorization\_descent}\),
\(\mathrm{k3\_to\_e\_specialization}\), \(\mathrm{transitions}\), and
\(\mathrm{scalar\_firewall}\); without them the compact source is only
a named input.  The finite packet
\(\texttt{certificates/e3\_source/k3e\_e3\_holfa}\) currently records
these rows only by the obstruction ledger
\(\texttt{E3\_SOURCE\_OBSTRUCTION\_LEDGER\_VERIFIED}\), with
\(\texttt{e3\_source\_certification=false}\).  The
boundary algebra \(A_b=\mathrm{End}_{\mathcal C}(b)\) on the open
factorization side can coincide with this target only after the
chain-fusion boundary vacuum \(b(K3\times E,\Sigma_2,E)\) exists.
Compact non-toric
\(K3\times E\) generalizations to other compact non-toric
\(\mathrm{CY}_3\)-data (e.g.\ Borcea--Voisin threefolds, Schoen's
small resolutions) require independent boundary vacua.

\paragraph{Boundary vacuum condition.} The missing object is
\(b(K3\times E,\Sigma_2,E)\), a boundary vacuum in an open
factorization dg-category on
\((E,D_E,\tau_E)\) with \(D_E\) at the Igusa cusp
\(\mathfrak c_\infty\) of the Siegel space.  The chain-fusion
isomorphism on \(W_{\le3}\) at finite HN height \(R=3\) must identify
the chiral algebra with
\(A_{b(K3\times E,\Sigma_2,E)}\). The necessary
compatibility condition is the consistency of the chain-fusion isomorphism with
the unscaled trace \(\Delta_5^{-2}\), whose OP representative is
\(-4096\,\Delta_5^{-2}\), and with the universal Borcherds-weight
identity \(\kappa_{\mathrm{BKM}}(\Phi_1)=c_1(0)/2=5\); see the
\(\kappa\)-tuple discussion below.

\subsubsection{The gravity-line operator algebra}
\label{subsubsec:open-gravity-line}

\paragraph{Gravity-line algebra.} A gravity-line operator algebra, if constructed,
is an operator algebra on the closed-colour factorization algebra
\(F^{\mathrm{cl}}\) on \(K3\times E\) whose Pentagon-face determinant
section is required to be \(\Delta_5^2=\Phi_{10}^{\mathrm{un}}\), a
weight-\(10\) automorphic section.  Its partition character is the
reciprocal section \(\Delta_5^{-2}\).  The primitive square-root row has
Borcherds weight
\[
\omega_{\mathrm{Borcherds}}(\Delta_5)
=\kappa_{\mathrm{BKM}}(\Delta_5)=c_1(0)/2=5.
\]
This is the same Hall--Borcherds residual and remains open in Vol~II.

\paragraph{Comparison with Vol II.} Vol~II \cite{ChiralBarCobarVolII}
names the gravity-line Hall--Borcherds comparison conjecture and the
Pentagon-face determinant;
Vol~III \cite{CYQGVolIII} names
\(\kappa_{\mathrm{BKM}}(\Phi_N)=c_N(0)/2\) as the universal
Borcherds-weight identity \cite{Borcherds95}, with
\(\kappa_{\mathrm{BKM}}(\Phi_1)=c_1(0)/2=5\) at the paramodular
\(\Delta_5\) datum.

\paragraph{The \(\kappa\)-tuple.} The five-fold \(\kappa\)-tuple
\((\kappa_{\mathrm{cat}},\kappa_{\mathrm{ch}}^{\mathrm{Hodge}},
\kappa_{\mathrm{ch}}^{\mathrm{Heis}},\kappa_{\mathrm{BKM}},\kappa_{\mathrm{fiber}})
=(0,0,3,5,24)\) at \(K3\times E\) records five separate invariants,
distinct in construction and meaning, never additively combined. Here
\(\kappa_{\mathrm{cat}}=\chi(\mathcal O_{K3\times E})=
\chi(\mathcal O_{K3})\cdot\chi(\mathcal O_E)=2\cdot 0=0\) is the
K{\"u}nneth-multiplicative categorical Euler;
\(\kappa_{\mathrm{ch}}^{\mathrm{Hodge}}=\sum_q(-1)^qh^{0,q}(K3\times E)=0\)
by Serre duality on the compact total space;
\(\kappa_{\mathrm{ch}}^{\mathrm{Heis}}=3\) is the rank of the
Heisenberg--Cartan in the chiral Hall specialization on
\(\mathrm{Ran}(E)\) and equals the Cartan rank of
\(\mathfrak g_{\Delta_5}\);
\(\kappa_{\mathrm{BKM}}=c_1(0)/2=5\) is the universal Borcherds-weight
identity at the Gritsenko--Cl{\'e}ry row whose Jacobi seed is
\(\phi_{0,1}\) and whose Borcherds product is \(\Delta_5\), with
\(c_1(0)=10\);
\(\kappa_{\mathrm{fiber}}=24\) is the Mukai-lattice rank
\(\mathrm{rk}\,\widetilde H^*(K3,\mathbb Z)=24\) of the K3 fibre, equal
to \(\chi_{\mathrm{top}}(K3)\). The additive identity
\(\kappa_{\mathrm{BKM}}=\kappa_{\mathrm{ch}}+\chi(\mathcal O_{\mathrm{fiber}})\)
fails because it confuses
\(\chi_{\mathrm{top}}(K3)=24\) with
\(\chi(\mathcal O_{K3})=2\); the universal identity
\(\kappa_{\mathrm{BKM}}(\Phi_N)=c_N(0)/2\) is the correct frame.

\paragraph{Argument.} The Pentagon-face determinant is the
modularity-witness of the open-side cyclic trace plus clutching;
modularity is a property of the open trace, not of the closed
algebra.  On the \(\Delta_5\) row, \(c_1(0)/2=5\) is the weight of the
primitive Borcherds factor.  The determinant required by chain fusion
is its square
\(\Delta_5^2=\Phi_{10}^{\mathrm{un}}\), hence a form of total weight \(10\)
on the closed cobordism; the level-\(n\) protected scalar shadow is
the reciprocal Laurent expansion \(\Delta_5^{-2}\).

\paragraph{Gravity-line trace condition.} The required object is the
gravity-line operator algebra
\(\mathrm{GravLine}_{K3\times E}\) on the open factorization
dg-category on \((E,D_E,\tau_E)\) at the Igusa cusp, with
Pentagon-face determinant identity
\(\det_{\mathrm{Pentagon}}^{\mathrm{cl}}(\mathrm{GravLine}_{K3\times E})
=\Delta_5^2\) and compatibility with the universal
Borcherds-weight identity
\(\kappa_{\mathrm{BKM}}(\Phi_1)=5\) and with the \(\kappa\)-tuple
\((0,0,3,5,24)\) at the K3-Heisenberg witness. The compatibility
condition is the existence of chain fusion at
\(b(K3\times E,\Sigma_2,E)\) and the Hall--Borcherds residual lifting
the protected trace
\(Z_{\mathrm{BPS}}^{K3\times E}=\Delta_5^{-2}\) toward the acting
operator-algebra level.

The finite compact-source recognition datum \((S1)\)--\((S10)\) on the
first relation-closed window is the packet
\[
\mathrm{window\_closure},\quad
\mathrm{target\_source\_representatives},\quad
\mathrm{chevalley\_serre\_matrices},\quad
\mathrm{hall\_boundary\_complex},\quad
\mathrm{spectral\_sequence},\quad
\mathrm{kernel\_matrices},\quad
\mathrm{pbw\_graded},\quad
\mathrm{first\_window\_theorems},\quad
\mathrm{transitions},\quad
\mathrm{scalar\_firewall}.
\]
The compact Hall source packet presently has a verified obstruction
ledger,
\[
\texttt{certificates/sources/k3e\_compact\_hall/blocked\_obligations.csv},
\]
which lists the missing \(M,D,\eta,\epsilon,B,G,K,Q,A\) matrices, the
Hall bialgebra identity rows, Hopf-pairing identity rows,
radical ideal/coideal rows, relation rows, no-extra equality,
generation, PBW, strict transition, \(A_\beta\)-intertwining, and
finite Koszul comparison rows.  Its verifier records
\(\texttt{compact\_source\_recognition=false}\).  This ledger is the
finite obstruction list for the source theorem, not a substitute for
the source theorem.
The first relation-closed window has its own obstruction ledger,
\[
\texttt{certificates/first\_window/k3e\_relation\_closed\_window/blocked\_obligations.csv},
\]
with status
\(\texttt{FIRST\_WINDOW\_OBSTRUCTION\_LEDGER\_VERIFIED}\) and
\(\texttt{first\_window\_certification=false}\).  It records the
missing window, representative, relation, Hall-boundary, spectral,
kernel, PBW, theorem, and transition rows for the first finite
primitive-recognition theorem.
The Dirac--Igusa pro-object morphism packet has its own obstruction
ledger,
\[
\texttt{certificates/morphisms/k3e\_dirac\_igusa\_pro\_object/blocked\_obligations.csv},
\]
with status
\(\texttt{MORPHISM\_OBSTRUCTION\_LEDGER\_VERIFIED}\) and
\(\texttt{pro\_object\_certification=false}\).  It records the missing
cofinal HN subsystem, fourteen finite-stage component rows, strict
transition morphisms, component maps for \((M1)\)--\((M8)\),
composition laws, Mittag--Leffler exactness, and pro-isomorphism rows.
The theorem-dependency ledger
\[
\mathrm{beilinson\_levels},\quad
\mathrm{objects},\quad
\mathrm{morphisms},\quad
\mathrm{theorem\_dependencies},\quad
\mathrm{dependency\_edges},\quad
\mathrm{forbidden\_edges},\quad
\mathrm{status\_separation},\quad
\mathrm{level\_equalities},\quad
\mathrm{equality\_sites},\quad
\mathrm{definition\_separation},\quad
\mathrm{obstruction\_independence},\quad
\mathrm{transitions},\quad
\mathrm{scalar\_firewall}
\]
is now a populated schema-complete ledger.  Its verifier checks
Beilinson-level sites, typed objects, morphisms, theorem rows, allowed
proof edges, forbidden proof edges, claim-status separation,
Beilinson-level equality scopes, equality-site rows, transition rows,
definition-separation rows, obstruction-independence rows, and
scalar-firewall coverage; it still records
\(\texttt{dependency\_certification=false}\) and
\(\texttt{mathematical\_certification=false}\).  Thus the proof-level
exclusions are finite rows, but the retained construction packets,
recognition matrices, mirror comparison, and gravity-line operator
algebra remain separate mathematical obligations.  Together with these
packets, the chain-fusion boundary vacuum
\(b(K3\times E,\Sigma_2,E)\), the Hall--Borcherds residual lifting
\(Z_{\mathrm{BPS}}^{K3\times E}=\Delta_5^{-2}\), the gravity-line
operator algebra with Pentagon-face determinant \(\Delta_5^2\), and the
Mathieu-twined denominators
\(\operatorname{den}(\mathfrak g_{\Delta_5}^{(g)})\) on the \(26\)
\(M_{24}\)-conjugacy classes are the unresolved data.  The four
obstructions concern, respectively, the D0-degeneration of the reduced
orientation theory, quotient orientation, Weyl transport of the
orientation line, and the type-II Pfaffian wall atlas.  The remaining
operator obstructions concern the Hall--Borcherds operator lift, compact
non-toric chain fusion, and the gravity-line algebra.  The Beilinson cut
separates these objects before any scalar trace is taken.
